\markboth{SYSTEM ARCHITECTURE}{SYSTEM ARCHITECTURE}
\section{System Architecture}\label{sec:architecture}

This section presents Training Copilot's three-layer architecture: the MCP data layer, the agent layer, and the LLM integration seam. The guiding principle is \textit{tool-agnostic extensibility}---no component names a specific tool or data source. Tools are discovered at runtime, servers are registered declaratively, and the LLM provider is resolved from configuration, so the stack can be extended without code changes.

\subsection{Design Principles}\label{sec:architecture:principles}

The architecture follows six principles derived from the MCP specification and microservice design patterns~\citep{dragoni_microservices_2017}: the model discovers tools via \texttt{list\_tools()} and decides autonomously which to call (\textit{tool-agnostic}); \texttt{ToolHost.call\_tool()} is the single entry point for every invocation (\textit{one call path}); each MCP server is a self-contained process with its own port, authentication, and lifecycle (\textit{servers as independent services}); unreachable servers are silently skipped rather than failing the system (\textit{discovery over hardcoding}); credentials are passed as connection headers, never exposed to the model (\textit{auth separated from declaration}); and switching LLM providers is a configuration change (\textit{vendor-neutral LLM seam}).

\subsection{Design Rationale}\label{sec:architecture:rationale}

Each major decision maps to a limitation identified in the state of the art (Section~\ref{sec:sota:summary}). Sports-science consensus statements~\citep{bourdon_monitoring_2017, kellmann_recovery_2018} establish that training decisions require data from multiple heterogeneous sources, and direct API integration would create exactly the $M \times N$ fragmentation \cite{yang_survey_2025} identified as the core scalability barrier for tool-augmented agents; MCP's standardised discovery and invocation reduces this to $M + N$, so each new source costs one server and one configuration line, with no changes to the host, agents, or UI. Five specialists rather than one follows from \cite{qu_tool_2024} and \cite{patil_gorilla_2023}, who showed that tool-selection accuracy degrades as the tool set grows---confirmed in our own early prototyping, where a single agent with all 40\,+ tools (excluding the optional 116-tool Telegram bridge) frequently picked the wrong one or conflated domains (e.g.\ using weather data to answer recovery questions). Scoping each specialist to at most three MCP servers restores accuracy and enables domain-specific prompts (the recovery specialist interprets HRV relative to baseline; the route specialist geocodes before routing), following the sports-science distinction between internal load, external load, environmental factors, route planning, and evidence-based guidance~\citep{bourdon_monitoring_2017, gabbett_training_2016, kellmann_recovery_2018}.

LangGraph~\citep{langchain_langgraph_2024} implements the specialists because it supports the ReAct loop~\citep{yao_react_2023} of interleaved reasoning and tool calling while letting the orchestrator treat each specialist as an opaque unit: it integrates natively with LangChain's tool abstraction (our MCP-to-LangChain bridge needs no adapter code), provides built-in MLflow autologging for the hierarchical span trees the evaluation framework depends on (Section~\ref{sec:evaluation:trace}), and gives each specialist a clear process boundary as an independent A2A server with its own scoped ToolHost---advantages neither AutoGen~\citep{wu_autogen_2023} (conversation-pattern-based, less suited to independent specialist processes) nor a custom prompt loop (no built-in state management or tracing) offers. A2A~\citep{google_a2a_2025} governs inter-agent communication because the orchestrator--specialist pattern needs capability advertisement, task delegation, and result collection without coupling agents to each other's internals: Agent Cards advertise capabilities, task-based delegation returns structured artifacts, and opaque internals let specialists be replaced or extended independently, while also letting agents run in-process, as separate local processes, or in separate containers with only a URL change.

\subsection{Architecture Overview}\label{sec:architecture:overview}

Figure~\ref{fig:architecture} shows the complete architecture. The data path flows from the UI through the agent layer to MCP servers and back; the chat path adds A2A coordination on top.

\begin{figure}[ht]
\centering
\resizebox{\textwidth}{!}{%
\begin{tikzpicture}[
    node distance=0.7cm and 0.8cm,
    every node/.style={font=\small},
    uibox/.style={draw=tcblue, fill=tcblue!8, rounded corners=3pt, minimum width=2.6cm, minimum height=0.7cm, thick, align=center},
    corebox/.style={draw=tcgreen, fill=tcgreen!8, rounded corners=3pt, minimum width=2.4cm, minimum height=0.7cm, thick, align=center},
    agentbox/.style={draw=tcyellow!80!black, fill=tcyellow!8, rounded corners=3pt, minimum width=2.0cm, minimum height=0.7cm, thick, align=center},
    serverbox/.style={draw=tccyan, fill=tccyan!6, rounded corners=3pt, minimum width=1.6cm, minimum height=0.7cm, thick, align=center},
    arr/.style={-{Stealth[length=5pt]}, thick},
    darr/.style={-{Stealth[length=5pt]}, thick, dashed},
    layerlabel/.style={font=\scriptsize\bfseries, rotate=90, anchor=south},
]

    % UI Layer (blue)
    \node[uibox] (ui1) {Streamlit\\Chat Tab};
    \node[uibox, right=0.6cm of ui1] (ui2) {React\\Frontend};
    \node[uibox, right=0.6cm of ui2] (ui3) {Telegram\\Bridge};

    % Core layer (green)
    \node[corebox, below=1.0cm of ui2] (orch_adapter) {Orchestrator\\Adapter};

    % Agent layer (golden) --- laid out as a single left-to-right chain so
    % adjacent boxes are always separated by an explicit gap, regardless of
    % how wide any individual label renders (avoids overlap bugs).
    \node[agentbox, below=1.0cm of orch_adapter] (orch_agent) {Orchestrator\\Agent :9000};
    \node[agentbox, below=1.0cm of orch_agent, xshift=-3.4cm] (load) {Load\\:9002};
    \node[agentbox, left=0.5cm of load] (recovery) {Recovery\\:9001};
    \node[agentbox, right=0.9cm of load] (context) {Context\\:9003};
    \node[agentbox, right=1.0cm of context] (route) {Route\\:9004};
    \node[agentbox, right=0.5cm of route] (fitness) {Fitness\\:9005};

    % MCP Server layer (cyan) --- one left-to-right chain with fixed gaps, so
    % adding a server never causes overlap (the enclosing resizebox scales the
    % whole row to \textwidth). Larger gaps separate the per-agent clusters.
    % Garmin+Strava+Flythrough are the Load cluster (Flythrough :8107 is scoped
    % to the Load specialist, see AGENT_MCP_SCOPE); Weather+Calendar to Context;
    % Routes+Google Maps to Route; Vector DB to Fitness.
    \node[serverbox, below=1.5cm of recovery] (garmin) {Garmin\\:8104};
    \node[serverbox, right=0.35cm of garmin] (strava) {Strava\\:8103};
    \node[serverbox, right=0.35cm of strava] (flythrough) {Flythrough\\:8107};
    \node[serverbox, right=0.7cm of flythrough] (weather) {Weather\\:8101};
    \node[serverbox, right=0.35cm of weather] (calendar) {Calendar\\:8105};
    \node[serverbox, right=0.7cm of calendar] (routes) {Routes\\:8102};
    \node[serverbox, right=0.35cm of routes] (gmaps) {Maps\\:8108};
    \node[serverbox, right=0.7cm of gmaps] (vectordb) {Vector DB\\(RAG)};

    % Layer labels --- all share one left edge (1.2cm left of the
    % left-most box, Recovery/Garmin's column), rather than each being
    % hung 1.2cm off its own nearest box, which misaligned them.
    \node[layerlabel, text=tcblue] at ([xshift=-1.2cm]recovery.west |- ui1) {UI Layer};
    \node[layerlabel, text=tcgreen] at ([xshift=-1.2cm]recovery.west |- orch_adapter) {Core};
    \node[layerlabel, text=tcyellow!80!black] at ([xshift=-1.2cm]recovery.west |- orch_agent) {Agents (A2A)};
    \node[layerlabel, text=tccyan] at ([xshift=-1.2cm]recovery.west |- garmin) {MCP / Data};

    % UI → Core
    \draw[arr, color=tcblue!60] (ui1.south) -- (orch_adapter.north west);
    \draw[arr, color=tcblue!60] (ui2.south) -- (orch_adapter.north);
    \draw[arr, color=tcblue!60] (ui3.south) -- (orch_adapter.north east);

    % Core → Orchestrator Agent
    \draw[arr, color=tcgreen] (orch_adapter.south) -- node[right, font=\scriptsize]{A2A} (orch_agent.north);

    % Orchestrator → Specialists (A2A)
    \draw[arr, color=tcyellow!80!black] (orch_agent.south) -- (recovery.north);
    \draw[arr, color=tcyellow!80!black] (orch_agent.south) -- (load.north);
    \draw[arr, color=tcyellow!80!black] (orch_agent.south) -- (context.north);
    \draw[arr, color=tcyellow!80!black] (orch_agent.south) -- (route.north);
    \draw[arr, color=tcyellow!80!black] (orch_agent.south) -- (fitness.north);

    % Specialists → MCP Servers (primary paths, solid): each specialist's
    % full MCP scope (core/config.py AGENT_MCP_SCOPE) gets a solid arrow,
    % since Context genuinely needs both Weather and Calendar for every
    % scheduling query, not one as a fallback of the other.
    \draw[arr, color=tccyan] (recovery.south) -- (garmin.north);
    \draw[arr, color=tccyan] (load.south) -- (strava.north);
    \draw[arr, color=tccyan] (load.south) -- (flythrough.north);
    \draw[arr, color=tccyan] (context.south) -- (weather.north);
    \draw[arr, color=tccyan] (context.south) -- (calendar.north);
    \draw[arr, color=tccyan] (route.south) -- (routes.north);
    \draw[arr, color=tccyan] (route.south) -- (gmaps.north);
    \draw[arr, color=tccyan] (fitness.south) -- (vectordb.north);

    % Load → Garmin (secondary/fallback path, dashed): Strava is the
    % specialist's primary activity source, Garmin only a fallback
    % (Section~4.5); edge points (not corners) keep the curve's landing
    % angle clean on both boxes.
    \draw[darr, color=tccyan!60] ($(load.south)+(-0.3,0)$) to[bend right=20] ($(garmin.north)+(0.35,0)$);

\end{tikzpicture}%
}% end resizebox
\caption[System architecture]{Training Copilot system architecture. Solid arrows indicate primary communication paths; dashed arrows indicate secondary/fallback paths. Agents communicate via A2A; specialists access data via scoped MCP connections.}\label{fig:architecture}
\end{figure}


\subsection{Layer 1: The MCP Data Layer}\label{sec:architecture:mcp}

The foundation is the Model Context Protocol~\citep{anthropic_mcp_2024}, a JSON-RPC\,2.0-based~\citep{jsonrpc_spec_2013} client--server protocol in which a \textit{host} discovers tools from one or more \textit{servers} and routes invocations transparently. Training Copilot uses the Streamable HTTP transport (stateless, one request per call), so each server is an independent HTTP endpoint requiring no persistent connection. The server registry is a declarative mapping from logical names to endpoints in one file (\texttt{core/config.py}, Listing~\ref{lst:registry} in the Appendix); each URL is read from an environment variable, falling back to a localhost default, so switching from local processes to Docker Compose containers is a URL change, not a code change.

\texttt{ToolHost} implements two operations. \texttt{list\_tools()} opens a fresh session to every registered server \textit{in parallel} via \texttt{asyncio.gather}, namespacing returned tools as \texttt{server\_\_tool} (dots are not legal in OpenAI function names); each tool's JSON-Schema parameters and docstring are preserved verbatim, since the docstring is the only text the LLM reads when deciding whether and how to call a tool. Servers that fail to connect within 60 seconds are silently skipped, returning an empty list rather than raising---the mechanism behind graceful degradation: the system functions with whatever servers are reachable. \texttt{call\_tool(name, args)} splits the namespaced name, routes to the right server, and returns the result as a JSON string; errors, including timeouts and connection failures, come back as structured \texttt{\{"error": "..."\}} objects rather than exceptions, so agents can reason about failures and report them rather than crash. The async implementation opens one fresh event loop per synchronous call via a bridge function, avoiding deadlocks when called from thread-pool workers such as the Streamlit runtime. Each server is a self-contained FastMCP service (Listing~\ref{lst:server_pattern} in the Appendix); per-server credentials are passed as connection headers via \texttt{ToolHost(headers=\{...\})}, keeping authentication separate from tool declaration so credentials never enter the model's context and tool definitions remain portable across deployments.

\subsection{Layer 2: The Agent Layer}\label{sec:architecture:agents}

The agent layer combines LangGraph for within-agent execution with A2A for between-agent coordination. Each of the five specialists is a LangGraph ReAct agent with access to a \textit{scoped} ToolHost that discovers tools only from its assigned MCP servers (Table~\ref{tab:agent_scope}), enforced at the infrastructure level---a specialist physically cannot call tools outside its domain, rather than relying on prompt instructions a model might ignore under complex queries.

\begin{table}[ht]
\begin{center}
\begin{minipage}{\textwidth}
\caption[Agent--MCP scope mapping]{Agent-to-MCP-server scope mapping. Each specialist discovers tools only from its assigned servers.}\label{tab:agent_scope}
\begin{tabularx}{\textwidth}{llX}
\toprule
\textbf{Agent} & \textbf{MCP Servers} & \textbf{Domain} \\
\midrule
Recovery (:9001) & Garmin & Sleep, HRV, Body Battery, stress, readiness \\
Load (:9002) & Strava, Garmin & Training load (CTL/ATL/TSB), volume, splits, zones \\
Context (:9003) & Weather, Calendar & Forecast, UV, pollen; calendar read/write \\
Route (:9004) & Routes & Route planning, geocoding, trail search, isochrones \\
Fitness (:9005) & \textit{(none---RAG)} & Exercise science from a vector DB of literature \\
\botrule
\end{tabularx}
\end{minipage}
\end{center}
\end{table}


The bridge between MCP tools and LangGraph agents is \texttt{core/mcp\_langchain.py}: \texttt{scoped\_host(server\_names)} constructs a \texttt{ToolHost} restricted to a subset of the registry, and \texttt{build\_tools(host, recorder)} wraps each discovered tool as a LangChain \texttt{StructuredTool} implementing the record-full/clip-for-model pattern (Section~\ref{sec:agents:recording})---the full JSON result is appended to a shared \texttt{recorder} list (becoming the A2A DataPart artifact), while the LLM receives a truncated copy where large arrays (GPS points, intraday timelines) are replaced with placeholders such as \texttt{[847 items]} and the total string is capped at 6{,}000 characters. The orchestrator (port 9000) is a LangGraph agent whose only tools are \texttt{ask\_<specialist>} functions, each a thin wrapper around \texttt{core/a2a\_client.call\_agent()}: it resolves the specialist's Agent Card, opens a streaming A2A session, and demuxes the response into status-update messages (relayed as progress to the UI), text artifact parts (the answer), and data artifact parts (the specialist's full tool-call records). The orchestrator has no MCP access of its own---it decomposes the query, delegates to specialists in parallel when domains are independent (the LLM emits multiple \texttt{ask\_*} calls in one step, executed concurrently by LangGraph), synthesises responses, and assembles the \texttt{trace} structure the UI uses for maps, charts, and the debug panel. Agents run non-streaming internally (\texttt{ainvoke}, not \texttt{astream}) for robustness against the KIT university gateway, which has shown instabilities under streaming load; progress instead surfaces through A2A status messages.

\subsection{Layer 3: The LLM Integration Seam}\label{sec:architecture:llm}

The LLM seam (\texttt{core/llm.py}) is a vendor-neutral abstraction over three backends: OpenAI-compatible endpoints (including the KIT gateway), official OpenAI, and Google Gemini via \texttt{ChatGoogleGenerativeAI}. It re-reads \texttt{.env} per call, so provider changes via the Settings UI (Figure~\ref{fig:screenshot_settings} in Appendix~\ref{sec:appendix:screenshots}) take effect without restarting agent processes; Table~\ref{tab:env} in Appendix~\ref{sec:appendix:env} lists the key environment variables that configure this and every other layer.

Figure~\ref{fig:ui_grid} gives a thumbnail overview of all seven Training Copilot UI tabs described throughout this section; Appendix~\ref{sec:appendix:screenshots} reproduces each at full size.

\begin{figure}[ht]
\centering
\newcommand{\uithumb}[2]{%
  \begin{minipage}{0.48\textwidth}
  \centering
  \small\textbf{#1}\\[2pt]
  \fbox{\parbox{0.95\linewidth}{\centering\vspace{1.1cm}\textcolor{tcgrey!40}{\small\textit{#2}}\vspace{1.1cm}}}
  \end{minipage}%
}
\uithumb{Chat}{Multi-domain chat + trace}\hfill
\uithumb{Health}{Sleep, HRV, Body Battery}

\vspace{10pt}

\uithumb{Dashboard}{Goal cockpit, load trends, activities}\hfill
\uithumb{Coach}{Race goal, milestones, zones, plan}

\vspace{10pt}

\uithumb{Flythrough}{3D GPS flythrough (Dashboard modal)}\hfill
\uithumb{Telegram}{Proactive coach + route delivery}

\vspace{10pt}

\uithumb{Settings}{OAuth, LLM, Telegram}

\caption[Training Copilot UI overview]{Overview of the Training Copilot UI across its five tabs, plus the 3D flythrough modal launched from the Dashboard and the Telegram delivery channel. Full-size screenshots with detailed descriptions appear in Appendix~\ref{sec:appendix:screenshots}.}\label{fig:ui_grid}
\end{figure}


\subsection{Supporting Subsystems}\label{sec:architecture:supporting}

Conversation history (\textit{per-user memory}) is embedded with the same local \texttt{all-MiniLM-L6-v2} model~\citep{reimers_sentence-bert_2019} as the fitness RAG pipeline and retrieved by cosine similarity to inject relevant prior exchanges into the prompt; a background distillation thread periodically condenses history into a human-readable profile (\texttt{soul.md}) of durable facts---training goals, injury history, schedule constraints---prepended to every subsequent prompt, with both layers degrading gracefully if unavailable. Two designed extensions build on this profile store (planned, not yet shipped in the evaluated prototype). Rather than waiting for durable facts to be inferred conversationally over several sessions, an \textit{onboarding questionnaire} on first login would ask directly for the target event, its date, a target pace or finishing time, and an optional free-text field for constraints such as injuries, writing the answers into the same durable profile so every specialist has them from the first conversation onward. \textit{Goal-plan tracking} would then compare the profile's target against live data: a lightweight scheduler re-invokes the orchestrator on a fixed cadence with a synthesised prompt weighing the current CTL/ATL/TSB trend and recovery status against the onboarding target, and---only once the deviation crosses a threshold, e.g.\ a missed long run or TSB drifting negative ahead of a taper week---pushes the resulting summary unprompted through the same outbound channel the Telegram bridge already provides (Section~\ref{sec:data:telegram}), rather than waiting for the user to ask, with a dashboard panel plotting plan against actual trend. Crucially, both reuse the orchestrator and existing specialists unchanged---only the trigger (schedule vs.\ chat message) and a small plan store would be new, so the agent count in Table~\ref{tab:agent_scope} is unaffected; Section~\ref{sec:agents:delegation} describes the resulting fourth delegation pattern. \textit{LLM-driven chart generation} has the model produce Plotly Python code executed server-side in a restricted namespace from a completed turn's tool results; a failed execution gets a single Reflexion-style retry~\citep{shinn_reflexion_2023} with the error as feedback, and still-failing charts are dropped, avoiding hardcoded chart types. \textit{Multi-channel delivery} exposes a uniform \texttt{run()} interface consumed by three frontends---a React app over FastAPI SSE, a Telegram userbot bridge with voice-memo transcription and portable route formats, and the original Streamlit chat tab---all sharing the same agent layer and differing only in rendering. The \textit{3D activity flythrough} is a cinematic camera animation along GPS tracks over satellite basemaps, with bearing, pitch, and zoom EMA-smoothed for fluidity, exported either in-browser via WebCodecs or server-side as MP4 through a headless Chromium path. It is launched two ways: directly from the Dashboard tab, and from the chat agent---the flythrough is exposed as its own single-tool MCP server (\texttt{prepare\_flythrough}, port 8107) scoped to the load specialist (Table~\ref{tab:agent_scope}), so once the user confirms orientation, map style, and duration the specialist packages a render request that the trace assembler lifts into an inline chat action (Section~\ref{sec:agents:recording}). The render \textit{request} thus travels the normal MCP path, but the map imagery itself does not: unlike every other external source in this paper, the flythrough's basemap and terrain tiles are not MCP-mediated---MapLibre GL, running in the browser, fetches free, keyless raster tiles directly (satellite imagery from ESRI's ArcGIS World Imagery service, elevation DEM tiles in Terrarium format from an AWS-hosted Open Data S3 bucket, and a CartoDB dark-matter vector style for the non-satellite mode), so this is the one visual layer Training Copilot renders without going through \texttt{ToolHost} at all (Table~\ref{tab:api_overview} in Appendix~\ref{sec:appendix:apis} lists these alongside every MCP-mediated source). Figure~\ref{fig:screenshot_flythrough} in Appendix~\ref{sec:appendix:screenshots} shows the resulting flythrough view.

\subsection{Deployment}\label{sec:architecture:deployment}

Training Copilot supports local development (each service a separate process started by \texttt{dev\_stack.sh}) and a containerised deployment via Docker Compose in which every MCP server and every agent runs in its own container, all built from a \textit{single shared Dockerfile}: MCP services select which server module to launch through a \texttt{SERVER} environment variable, agent services override the container command, and inter-service URLs are supplied as environment variables. Because the registry is declarative, switching between the local-process and container modes (or a hybrid of the two) is a URL change, not a code change.
